% right above this is label: sec:spv-in-ut

\subsubsection{Introduction}

Our goal in this section is to find a method for cross-chain transactions between any two simplex-chains that is safe, fast, succinct, and reliable.
\begin{itemize}
    \item \emph{Safe}: an adversary must perform a 51\% attack to execute an invalid cross-chain transaction.
    \item \emph{Fast}: $O(t) \le O(\lambda) = O(1)$, where: $t$ is the maximum delay between a cross-chain transaction being confirmed by a simplex-chain and the point at which the destination simplex-chain allows processing that transaction; $\lambda$ is a security parameter (e.g., $\lambda = 4$ local confirmations).
    \item \emph{Succinct}: for any cross-chain proof $\pi$, $O(\pi) < O(c)$ and $|\pi|$ is practical.
    \item \emph{Reliable}: the above properties hold for all transactions in all simplex-chains.
\end{itemize}

Given everything we have discussed up to now, it's not clear how cross-chain transactions are possible.
The problem is that miners on one simplex-chain do not validate the state transitions of other simplex-chains.
Each simplex-chain, locally, has an $O(\nicefrac{n}{c})$ share of the network's security (assuming there are $O(c)$ many chains).
An $O(\nicefrac{n}{c})$ adversary could theoretically produce \emph{all} of the blocks for a single simplex-chain (e.g., by mining at a slight loss to push other miners out).
The adversary could then create a very long chain of blocks that have valid headers, PoRs, etc, but also contain one or more fraudulent transactions.
Perhaps the main chain eventually corrects to a non-fraudulent history (similar to \autoref{sec:preventing-dos-attacks}) --- but in that case we'd need to wait a \emph{very} long time to be sure that such a correction \emph{must} have happened.
That violates the \emph{fast} goal, so this eventual-correction idea fails to meet our overall goal.

Are we stuck?
For a blockchain network to be scalable, it must be the case that miners are \emph{not} required to validate more than one chain!
Have we come all this way, attempting to deal with the core conflict of the trilemma, only to find that we just moved the conflict?

\subsubsection{Why does SPV work for Bitcoin?}

For the sake of simplicity, let's consider Bitcoin in isolation.
Particularly, we'll assume that there are no other networks of similar size using the same PoW algorithm or compatible mining hardware.

A concise argument for the security of Bitcoin-SPV proofs is:

\begin{argumentslist}
    \item An SPV proof is a merkle branch proving that a transaction exists in a block.
    \item All blocks are part of the same blockchain and history.
    \item All transactions that occur are in some block's merkle tree.
    \item If a transaction is invalid, then the block is invalid.
    \item If a block builds on an invalid block, then it is invalid.
    \item If a miner produces an invalid block, they get no block reward.
    \item $\implies$ Honest miners won't produce or build on invalid blocks.
    \item \label{prem:spv-7} $\implies$ For an adversary to produce an apparently valid (but fraudulent) proof leading to a \emph{main chain} block, they'd need to out-compete honest miners (i.e., 51\% attack the network).
    \item \label{prem:spv-8} $\implies$ Only valid transactions exist in the main chain.\footnote{
        This isn't the full picture for some blockchains (e.g., Ethereum), even though SPV still works for them --- a transaction that fails to execute can still be valid and produce a valid state transition.
        In these cases, an invalid state transition implies an invalid block, so it's customary to prove that some state has the right values rather than that a transaction exists.
    }
    \item $\implies$ It's safe to accept SPV proofs for blocks in the main chain.
\end{argumentslist}

It appears that \cref{prem:spv-7} is our current sticking point for Simplex-SPV proofs.

But, we don't necessarily \emph{need} to replicate \cref{prem:spv-7} if we could get to \cref{prem:spv-8} by other means.

\subsubsection{Method}

\infoNote{
    The details below are an overview of Amaroo's cross-chain protocol.
    Full details comprise the Amaroo team's forthcoming paper on the subject.
}

\subsubsubsection{Context}

The purpose of UT's cross-chain protocol is to \emph{maintain coherence} between simplex-chains.
This means that all the essential properties we'd expect to hold for transactions on a single chain should hold between chains, too: coins are not created nor destroyed\footnote{
    That isn't to say burning coins isn't possible, but it should never be an unintended byproduct of a cross-chain transaction.
}, coins can only be spent once, \emph{the transaction executes entirely or not at all}, etc.
There will obviously be some particular properties that are different, such as a delay between inputs being spent (on the sending chain) and outputs being created (on the receiving chain).



\begin{comment}
    - hard things
      - guaranteed execution vs guaranteed to be redeemable

    - if a tx is valid, then spending the same inputs with a cross-chain wrapper around the outputs is valid
    - the validity of a transaction is evaluated and enforced entirely on the sending chain
    - if a transaction is invalid, then a fraud proof is possible
    - the \AxiomOfAvailability means miners on the receiving chain can easily find all incoming XC transactions
    - if a fraud proof exists for a block, then the block is invalid
    - a fraud proof can be included by any miner from any simplex chain (and goes in the PoR graph)
    - if a block is invalid, then all blocks of that chain that build on it are invalid
    - if a block is invalid, reflecting blocks are not immediately invalidated
    - however, reflecting blocks outside the security parameter are invalid \emph{unless} they build on PoR graph tips that include the fraud proof


    - normal SPV is PUSH-PULL: you PUSH on the sending chain and PULL on the receiving chain
    - We can do PUSH-PUSH: you PUSH on the sending chain it is automatically PUSHed on the receiving chain
    - We can also do PULL-PUSH: you PULL on the receiving chain, which automatically PUSHes on the sending chain to the receiving chain
      - this requires an output type: SpendIfExistsElse which tests whether an output exists and then produces one of two outcomes depending on the result
      - you send that PUSH-PUSH style from a chain, which then auto-execs the output on the receiving chain. If the exists branch is an XC output, then this is also PUSH-PUSH. So it functions as PULL-PUSH.
    - Our cases require only 1 tx compared to the typical 2 txs.



- DoS issue
- Bribe issue
- Decoherence (mint funds)

% important: censorship resistance
\end{comment}


% -- context

Before we address how cross-chain transactions will work, we first need to remember the context in which we're operating.
Leading up to this point, we've constructed a network with some wildly different properties to traditional blockchains: ultra-fast confirmations, additional DoS and censorship resistance, shared security, higher order scaling, etc.
These improvements are due to the specific combination of: \emph{Proof of Reflection}, the use of chain-like block-DAGs in the PoR Graph and each simplex-chain, and the Axioms of \AxOfAvailability{}, \AxOfMaximalReflection{}, and \AxOfUnifiedAncestry{}.
Each of these ideas, along with fraud proofs, will have a role to play in the cross-chain protocol described shortly.

Additionally, we are concerned about some possible attacks aimed at the cross-chain protocol.
We are not necessarily concerned about how the protocol resolves during a 51\% attack in this discussion since we already know that the network does not function in that context.

The first and most important thing to consider is \emph{Under which circumstances can cross-chain transactions be invalidated?}
If it is possible to invalidate a past cross-chain transaction, this could be used to perform a doublespend, or coin duplication, etc.
Therefore, we should strive for a protocol that only permits such invalidation under a 51\% attack (or worse).
That way, we can be confident that the cross-chain protocol is not the weakest link in the chain, since local (intra-chain) transactions break down at that point, anyway.

The obvious way to invalidate a cross-chain transaction is to cause a reorganization on the origin chain.
If the pivot of the reorganization is antecedent to the cross-chain transaction, then it is up to the attacker whether that transaction is included in the resulting canonical history.
This outcome is self-evidently insecure.
Therefore our question about cross-chain transaction invalidation has the same answer as questions like \emph{Under which circumstances can a simplex-chain reorganization occur?}
We are also concerned with the parameters around such a reorganization: \emph{Is it dependent on the number of confirmations?}, \emph{If so, how many?}, \emph{What proportion of global hash-rate is required?}, etc.

Fortunately, we've done a lot of the legwork already to prove that the ordering of simplex-chain blocks is stable and works for cross-chain transactions.
\emph{Proof of Reflection} allows us to share security between simplex-chains.
Block-DAGs and the PoR Graph provide some censorship resistance, and more importantly, are the basis for both the \AxiomOfMaximalReflection{} and the \AxiomOfUnifiedAncestry{}.
These axioms go on to, respectively, bind the histories of each simplex-chain to one another and ensure that each chain's history spans all valid blocks of that chain in the PoR graph.
The binding of histories is the basis for the \emph{coherence} of the simplex, and PoR graph-spanning histories ensures that the ordering of blocks within each chain conforms to the ordering of blocks in the PoR graph.
Therefore: to change the order of blocks within a simplex-chain requires that an attacker change the global ordering, which in turn requires a 51\% attack.
$\blacksquare$

Regarding network security, resisting this particular attack is paramount because this kind of reorganization \emph{affects honest nodes}.
By comparison, there are other attacks that do not affect honest nodes, but do affect, say, light clients.\footnote{
    In this context light clients are not considered honest nodes (nor attacking nodes) since they do not contribute to global consensus.
}
These other attacks aren't as concerning provided that they only affect non-validating network participants.
But! What happens when an invalid block is mined?
It \emph{will} be reflected, because miners of other simplex-chains do not validate the block.
What we need is a method of \emph{error correction} -- one that is fast enough to ensure invalid blocks are dealt with before the cross-chain protocol allows those blocks to be processed.
This role is filled by \emph{fraud proofs}.

In \UT{}, \emph{fraud proofs} are automatically constructed by any full node during normal validation -- the state transition function will return either a new state trie \emph{or} a fraud proof.
Once a fraud proof is constructed, it is broadcast similar to a transaction.
Then, \emph{any miner} from \emph{any simplex-chain} can validate and include the fraud proof \emph{as part of the PoR graph}.
Due to \AxOfUnifiedAncestry{}, every simplex-chain calculates the expected chain-tips for every simplex-chain.
When a fraud proof invalidates a block, the calculated chain-tips for a given simplex-chain are updated to exclude any invalid blocks from that chain's history.
Thus, as long as fraud proofs are readily generated, broadcast, and included by other miners, the PoR graph knows which blocks are on-main valid blocks, and which are not.

There is, of course, some lag between an invalid block being mined, the fraud proof being generated, and the proof being included in the PoR graph.
In a healthy network, the expected delay is $(N_1 B_f)^{-1} + o$ seconds, where $o$ is a little overhead to account for the first node of that chain to receive and process the block.
This is much less than a block period -- even when a minority of miners are honest (particularly when $N_1 \gg p^{-1}$).

\subsubsubsection{Construction}

Typically, a cross-chain protocol will require two user interactions: one on the sending chain and one on the receiving chain.
That is, there is some state modification on the first chain, and then, on the second, an acknowledgment of that and some associated action (like crediting an account).
Based on the direction of information, order of actions, and the focus on those actions, we could describe these kinds of methods as \emph{push-pull} methods.
Information is `pushed' from the sending chain, and, at some later point, `pulled' into the receiving one -- conceptually, at least.
In practice the `pull' requires the user to publish a proof about the sending chain to the receiving chain.

By contrast, a \emph{push-push} method would be one where the receiving chain automatically processes the cross-chain output at the correct time.
This means that no second action from the user is required.
Such a system must therefore \emph{require} miners to process cross-chain transactions -- but how do they know about transactions on other chains and whether they are valid?
In our case, the \AxiomOfAvailability{} means that all miners on all simplex-chains have all recent blocks from all other simplex-chains.
Our use of fraud proofs means that all honest nodes quickly learn of invalid blocks.
Thus, provided enough time has passed, we can be confident that miners know of all  cross-chain transactions, and that they are, in fact, valid.

How long should we wait, and how do we measure how long we've waited?
We need at least some delay to avoid a particularly lucky series of blocks from triggering cross-chain processing.
Ideally, this processing should be as predictable as possible, and well distributed.
We also want to avoid our delay method from being manipulated to cause earlier or later execution.

The delay duration itself is almost unconstrained: fraud proofs both propagate through the network well within one block period (at which point honest nodes are aware of them), and are included in the PoR graph within a block period, too.
However, if we pick too short a period, then we risk some fragility in the network.
We can technically recover if an invalid cross-chain transaction was processed before a corresponding fraud proof was recorded, but doing so robs honest miners of their block reward.
So, whatever we pick, it needs to be something that is conservative enough for miners to be comfortable with.

Although it's somewhat arbitrary, 60 seconds (or 4 block periods at $B_f = \nicefrac{1}{15}$ Hz) seems like a reasonable choice.
There is plenty of time for the fraud proof to propagate and be recorded, and the risk of an honest miner's block being invalidated is low.

How do we measure this duration, then?
We could use block timestamps, although that may open us up to some kind of time-warp attack.
We could use the number of confirmations on the sending or receiving chain, but there's local block variance to consider.
We could use the change in volume of the simplex (i.e., the number of confirmations over all simplex chains since some block), but we don't have any guarantees about the progression of either the sending or receiving chain.
Each of these has disadvantages, but if we take these breakpoints together, then we have a reliable single condition.

We can now summarize; when drafting an \cL block, the \cR blocks which are valid for cross-chain processing from the draft block's perspective are those which:
\begin{enumerate}
    \def\labelenumi{\arabic{enumi}.}
    \tightlist
    \item Have not been processed before; and
    \item Have a timestamp at least $\lambda$ block periods in the past; and
    \item Have a volume at least $N_1 \lambda$ blocks fewer than this block's volume; and
    \item \label{e-xc-cond3} Are in the PoRs history of this block's $\lambda$-parent\footnote{
        Definition: the $1$-parent is the parent; the $2$-parent is the parent of the $1$-parent; the $(k+1)$-parent is the parent of the $k$-parent.
    }; and
    \item \label{e-xc-cond4} Are at least $\lambda$ local confirmations deep.
\end{enumerate}

These conditions are illustrated in \autoref{fig:xc-conditions}.

\begin{figure}
    \centering
    \includegraphics[max width=\linewidth]{spv/xc_security_conditions_sag}
    \caption[
        % The conditions for when a block becomes valid for cross-chain transaction processing.
        The conditions used to find \cR blocks valid for cross-chain transaction processing.
    ]{
        The conditions used to find \cR blocks that are valid for cross-chain transaction processing from the perspective of $L_{i}$.
        Particularly, we find up to 4 distinct blocks via 4 different paths.
        Note that these blocks can appear in any chronological order, and the specific blocks shown in the diagram are just one example of many possibilities.
        In this case,
        $\lambda$ block periods ago resolves to $\cR_{j-x_3}$;
        volume at least $N_1 \lambda$ blocks fewer resolves to $\cR_{j-x_2}$;
        the $\lambda$-parent of $R_j$ resolves to $R_{j-\lambda}$; and
        the $\lambda$-parent of $L_i$ resolves to $L_{i-\lambda}$.
    }
\label{fig:xc-conditions}
\end{figure}

\begin{comment}
    Visual representation of the conditions used to compute, for a block $L_i$ on chain $L$, blocks from chain $R$ that are valid for cross-chain transactions.
        Valid blocks are those that are earlier than (or in) a common ancestor of the four blocks $\{ R_{j-\lambda}, R_{j-x_1}, R_{j-x_2}, R_{j-x_3} \}$.
        $R_j$ is a PoR ancestor of $L_i$ on chain $R$.
        $R_{j-\lambda}$ is a $\lambda$-parent $R_j$.
        $L_{i-\lambda}$ is a $\lambda$-parent $L_i$.
        $R_{j-x_1}$ is a PoR ancestor of $L_{i-\lambda}$ on chain $R$.
        $R_{j-x_2}$ is an ancestor of $L_i$ such that the PoR volume has increased at least $N_1 \lambda$ since $R_{j-x_2}$.
        $R_{j-x_3}$ is an ancestor of $L_i$ that is at least ${B_f}^{-1} \lambda$ older than $L_i$.
\end{comment}

There is one slight oversight in this formulation that we need to fix.
How do we know that the transactions in those blocks to be processed are all consistent?
If we naively collect blocks according to the above conditions, then we don't -- there might be conflicting transactions in off-main blocks.
To correct this, we first use the four conditions to find the 1-4 most recent corresponding blocks (some conditions may resolve to the same block).
Next, for chain \cR, we find the on-main MRCA (most recent common ancestor) of those blocks -- this is the most recent block that will be processed for cross-chain transactions.
We repeat this process for the parent of the draft block to find its corresponding MRCA by the same rules.
Let's call these blocks $\cR_\alpha$ and $\cR_\beta$ respectively.
The set of blocks to process is: $\cR_\alpha$ and all on-main blocks in its history \emph{minus} $\cR_\beta$ and all on-main blocks in its history.
This is the on-main ``chain-slice'' from $\cR_\alpha$ (inclusive) to $\cR_\beta$ (exclusive).
Since all of these blocks are on-main, any valid transactions from off-main blocks between $\cR_\alpha$ and $\cR_\beta$ will be included in one of the on-main blocks.
By inspection, we can observe that $\cR_\beta$ for the current block was $\cR_\alpha$ for the parent block (or an earlier ancestor).
By induction, this series of chain-slices continues all the way back to the genesis block (before which no chain-slices exist).
Therefore, every on-main block (and thus every valid transaction on that chain) is processed for cross-chain transactions at the earliest appropriate time.

The protocol for processing cross-chain transactions is now, simply: when drafting an \cL block, for each \cR chain, miners gather the chain-slice (of \cR blocks) that is valid for cross-chain processing, scan those blocks for incoming cross-chain outputs, and include those cross-chain outputs in the draft block.

\aside{
    In practice, these conditions aren't ideal because updating the reflections of a draft block may introduce new blocks which are valid for cross-chain processing, in turn requiring a complete recalculation of that block's output state.
    To avoid this, we replace \autoref{e-xc-cond3} and \autoref{e-xc-cond4} with similar conditions based on the draft block's parent and $\lambda-1$.
    This keeps the set of blocks stable when updating the reflections of a draft block.
}

\subsubsubsection{Evaluation}

\subsubparagraph{Safe}

Double-spending a cross-chain transaction requires a 51\% attack on the entire network.
Executing an invalid cross-chain transaction is not possible when there are honest nodes of that simplex-chain and honest miners of any simplex chain.

\subsubparagraph{Fast}

Cross-chain transactions are processed based on conditions which are all $O(1)$ with respect to time, and therefore the expected delay in execution is $O(1)$.

\subsubparagraph{Succinct}

No proof is explicitly recorded.
If one did want to craft an explicit proof $\pi$, then it would be similar in size to a normal transaction proof (assuming headers from the surrounding PoR graph are available).
Normal transaction proofs are $O(\log c)$, therefore, $O(\pi) = O(\log c) < O(c)$.

\subsubparagraph{Reliable}

The consensus protocol requires that cross-chain transactions are processed if they are valid for cross-chain processing.
Not doing so invalidates a block.
Therefore, all cross-chain transactions on all simplex-chains will have these above properties.

\subsubparagraph{Censorship Resistance}

This construction has a curious property.
Say that an attacker, in an effort to perform an empty block DoS, is mining a simplex-chain at a loss, and has pushed out other miners.
If you wish to make a local transaction, it could take much longer than a block period.
An alternative is to make a cross-chain transaction to that simplex-chain from another simplex-chain instead.
This will obviously take a minute or so to execute, but you will have a guarantee that it will be processed -- the attacker cannot continue the attack without processing that cross-chain transaction.
You could then `rescue' your coins and send them to a more useful chain.\footnote{
    This kind of operation would require some support from the transaction layer, e.g., a way to spend inputs indirectly.
    Such an output type would also need to contain instructions in case of a failure, since the sending chain cannot verify that the inputs to be spent (on the receiving chain) will be spendable when the cross-chain transaction is processed.
    Since we must scan blocks as part of the cross-chain processing, the branch taken should also be recorded in the block, alongside the corresponding output or transaction.
}
While this describes an unpleasant user experience, the takeaway is that it is yet another obstacle for an attacker to overcome.


% hidden assumption: attacker is producing valid blocks.
% this matters particularly because \emph{the re-org affects honest nodes}.


% -- ctx: attacker is producing valid blocks
% problem: conflates DoS with XC attack
\begin{comment}
\subsubsubsection{Other draft stuff}
    \emph{The only re-orgs that matter are those among \ul{honest} nodes.}
    For example: if an attacker bribes 90\% of miners to ignore certain blocks, this \emph{does not affect \ul{\UT{1}}}.
    It \emph{does} affect traditional blockchains, though, since this \emph{effectively} controls the canonical history of the chain.
    For a single-chain network, using a block-DAG is enough to foil the attack: controlling the canonical history of the chain can only be done for very short intervals (\autoref{sec:preventing-dos-attacks}).
    In a multi-chain network like \UT{1}, we need the property to be strong enough to ensure that cross-chain transactions are \emph{never} processed when the network is in an ambiguous state.
\end{comment}
% \UT{1} achieves this through the \emph{convergence} required by the Axioms.
% \UT{1} achieves this

% one way to achieve this is by making certain things invalid that would be permissable (or not possible to detect) in a single-chain netowrk.
% UT does this via Axioms


\begin{comment}
- any honest node can produce a fraud proof
- any miner on any simplex-chain can include a fraud proof at the PoR level targeting any other chain
- therefore, getting invalid blocks in the main chain requires either:
  - 51% attack, or
  - near-complete control over p2p connections (particularly around miners)
\end{comment}

\subsubsection{Decoupled State Progression}

The key to cross-chain transactions working is the decoupling of state progression combined with the property that cross-chain transactions can be read from \cR blocks without needing to validate the \cR chain -- if it exists, then it is valid.
Since part of this protocol involves asynchronous communication between chains, we need a method of error correction to ensure cross-chain transactions are safe: fraud proofs.
Crucially, we expect that honest nodes (particularly those on other chains) are never fooled by an invalid block for long --- much less than a block period.
On this basis we can pick a safety parameter $\lambda$ which (in combination with the conditions from earlier) provides enough buffer to ensure that the PoR graph remains coherent.

This decoupling is the essential component of \UT{} that makes $O(c^2)$ scaling possible, and is a direct result of how we use PoR to create the PoR graph.
The procedural elements are shown in \autoref{fig:spv-network-complexity-split}.


\begin{figure}[htbp]
    \centering
    \includegraphics[max width=\linewidth]{spv/network_scale_AL_split_sag}
    \caption[
        Flow chart of how the segmentation of state works in \UT{} to achieve $O(c^2)$ capacity.
    ]{
        Flow chart of how the segmentation of state works in \UT{} to achieve $O(c^2)$ capacity.
        ``AL'' is short for ``Action Layer'', a term invented here so that we can point to the various stages of PoR graph extension and state progression.
        $\text{AL}^1$ columns indicate actions that validating nodes of any simplex-chain perform on the PoR graph.
        The $\text{AL}^2$ column indicates that only validating nodes \emph{of that chain} perform the associated action.
    }
    \label{fig:spv-network-complexity-split}
\end{figure}
